\documentclass[11pt,a4paper]{article}

% --- Encoding ---
\usepackage[utf8]{inputenc}
\usepackage[T1]{fontenc}
\usepackage[spanish,es-noshorthands]{babel}
\usepackage{lmodern}

% --- Layout ---
\usepackage[margin=2.5cm]{geometry}
\usepackage{parskip}
\setlength{\parindent}{0pt}
\setlength{\parskip}{0.5em plus 0.2em minus 0.1em}

% --- Tables and graphics ---
\usepackage{booktabs}
\usepackage{array}
\usepackage{enumitem}
\usepackage{xcolor}
\usepackage{graphicx}
\usepackage{tikz}
\usetikzlibrary{positioning,shapes.geometric,arrows.meta,fit,backgrounds,calc,decorations.pathreplacing}
\usepackage{caption}
\usepackage{subcaption}
\usepackage{amsmath,amssymb}

% --- Soporte para caracteres Unicode usados en formulas en texto plano ---
\usepackage{newunicodechar}
\newunicodechar{→}{\ensuremath{\rightarrow}}
\newunicodechar{≥}{\ensuremath{\geq}}
\newunicodechar{≤}{\ensuremath{\leq}}
\newunicodechar{×}{\ensuremath{\times}}
\newunicodechar{÷}{\ensuremath{\div}}
\newunicodechar{≈}{\ensuremath{\approx}}
\newunicodechar{⌈}{\ensuremath{\lceil}}
\newunicodechar{⌉}{\ensuremath{\rceil}}
\newunicodechar{✓}{\ensuremath{\checkmark}}

% --- Code listings ---
\usepackage{listings}
\definecolor{codegray}{rgb}{0.96,0.96,0.96}
\definecolor{codeborder}{rgb}{0.85,0.85,0.85}
\definecolor{codecomment}{rgb}{0.30,0.55,0.30}
\definecolor{codekeyword}{rgb}{0.10,0.30,0.70}
\definecolor{codestring}{rgb}{0.65,0.15,0.10}

\lstdefinestyle{cstyle}{
    language=C,
    backgroundcolor=\color{codegray},
    basicstyle=\ttfamily\footnotesize,
    breaklines=true,
    frame=single,
    rulecolor=\color{codeborder},
    showstringspaces=false,
    keywordstyle=\color{codekeyword}\bfseries,
    commentstyle=\color{codecomment}\itshape,
    stringstyle=\color{codestring},
    aboveskip=1em,
    belowskip=1em,
}

\lstdefinestyle{shellstyle}{
    language=bash,
    backgroundcolor=\color{codegray},
    basicstyle=\ttfamily\small,
    breaklines=true,
    frame=single,
    rulecolor=\color{codeborder},
    showstringspaces=false,
    keywordstyle=\color{codekeyword}\bfseries,
    commentstyle=\color{codecomment}\itshape,
    stringstyle=\color{codestring},
    aboveskip=1em,
    belowskip=1em,
    columns=fullflexible,
    upquote=true,
}

\lstdefinestyle{textstyle}{
    backgroundcolor=\color{codegray},
    basicstyle=\ttfamily\footnotesize,
    breaklines=true,
    frame=single,
    rulecolor=\color{codeborder},
    showstringspaces=false,
    aboveskip=1em,
    belowskip=1em,
    columns=fullflexible,
}

\lstset{style=textstyle}

% --- Definition / Idea / Important boxes ---
\usepackage[most]{tcolorbox}
\newtcolorbox{idea}[1][]{
    colback=orange!5!white,
    colframe=orange!60!black,
    fonttitle=\bfseries,
    title=Idea intuitiva,
    sharp corners=southwest,
    rounded corners=northwest,
    arc=2pt,
    boxrule=0.6pt,
    left=8pt,right=8pt,top=4pt,bottom=4pt,
    #1
}
\newtcolorbox{importante}[1][]{
    colback=red!5!white,
    colframe=red!50!black,
    fonttitle=\bfseries,
    title=Importante,
    sharp corners=southwest,
    rounded corners=northwest,
    arc=2pt,
    boxrule=0.6pt,
    left=8pt,right=8pt,top=4pt,bottom=4pt,
    #1
}
\newtcolorbox{respuesta}[1][]{
    colback=green!5!white,
    colframe=green!40!black,
    fonttitle=\bfseries,
    title=Respuesta,
    sharp corners=southwest,
    rounded corners=northwest,
    arc=2pt,
    boxrule=0.6pt,
    left=8pt,right=8pt,top=4pt,bottom=4pt,
    #1
}

% --- Hyperlinks ---
\usepackage[colorlinks=true,
            linkcolor=blue!60!black,
            urlcolor=blue!50!black,
            citecolor=blue!60!black]{hyperref}

% --- Color palette ---
\definecolor{c1}{HTML}{4C72B0}
\definecolor{c2}{HTML}{DD8452}
\definecolor{c3}{HTML}{55A868}
\definecolor{c4}{HTML}{8172B2}
\definecolor{cred}{HTML}{C44E52}
\definecolor{cfree}{HTML}{8C8C8C}

% --- Title block ---
\title{
    \vspace{-2.5em}
    \textbf{Examen Parcial 1 --- 2023} \\[0.3em]
    \large Sistemas Operativos \\
    \small Licenciatura en Ciencias de la Computación
}
\author{
    Tomás Rodeghiero \\
    \small UNRC --- Universidad Nacional de Río Cuarto
}
\date{2023}

\begin{document}

\maketitle

\begin{abstract}
\noindent
Resolución completa del Parcial 1 del año 2023 de la cátedra de
Sistemas Operativos (UNRC, Lic.\ en Ciencias de la Computación). El
parcial consta de diez ejercicios que recorren los temas fundamentales
del primer cuatrimestre: seguridad y modos de privilegio,
contraste entre kernel monolítico y microkernel, instrucciones
privilegiadas, mecánica de las syscalls, secuencia de
\texttt{fork}/\texttt{exec}/\texttt{wait} en un shell, bloqueo del
proceso ante una syscall lenta como \texttt{read}, transición entre
estados (\texttt{wakeup}), \emph{starvation} en distintos algoritmos
de planificación, \emph{race conditions} en accesos a un saldo
bancario compartido, y las limitaciones de spinlocks en
monoprocesador y de la inhibición de interrupciones en
multiprocesador. Cada ejercicio se desarrolla con la estructura
\textbf{Enunciado $\to$ Análisis $\to$ Desarrollo $\to$ Respuesta}.
\end{abstract}

\tableofcontents
\bigskip

% =====================================================================
\section{Ejercicio 1 --- Seguridad y modos de privilegio}
% =====================================================================

\textbf{Enunciado.} Un sistema de multi-programación o multi-tarea
debe compartir recursos.

\textbf{(a)} ¿Qué problemas de seguridad surgen?

\textbf{(b)} ¿Es posible lograr un sistema seguro con una CPU con un
solo modo de privilegio?

\subsection*{1.a --- Problemas de seguridad al compartir recursos}

Cuando varios procesos comparten CPU, memoria, dispositivos y
archivos, surgen riesgos clásicos:

\begin{itemize}[leftmargin=*]
    \item \textbf{Confidencialidad de la memoria}: sin protección,
        un proceso podría leer la memoria de otro y obtener datos
        sensibles (contraseñas, claves, datos privados).
    \item \textbf{Integridad de la memoria}: un proceso podría
        \emph{modificar} la memoria de otro (incluido el código del
        kernel), corrompiéndolo o inyectando código malicioso.
    \item \textbf{Acceso no autorizado a dispositivos}: lectura
        directa del disco/red bypaseando los permisos del FS;
        manipulación del audio o pantalla de otro usuario; etc.
    \item \textbf{Denegación de servicio (DoS)}: un proceso podría
        monopolizar la CPU (no ceder voluntariamente), acaparar toda
        la RAM, llenar el FS, o bloquear dispositivos compartidos.
    \item \textbf{Acceso no autorizado a archivos}: lectura/escritura
        de archivos de otros usuarios sin permiso.
    \item \textbf{Ejecución de instrucciones críticas}: si un proceso
        pudiera ejecutar instrucciones como ``modificar tabla de
        páginas'' o ``deshabilitar interrupciones'', podría
        secuestrar el sistema entero.
    \item \textbf{Escalada de privilegios}: un proceso de usuario
        común podría hacerse pasar por root o por otro usuario.
\end{itemize}

\subsection*{1.b --- ¿Sistema seguro con un solo modo de privilegio?}

\begin{respuesta}
\textbf{No}, en general. Sin distinción entre \emph{modo usuario} y
\emph{modo kernel}, cualquier programa puede ejecutar todas las
instrucciones de la CPU, incluidas las críticas:

\begin{itemize}[leftmargin=*,nosep]
    \item Modificar la tabla de páginas $\Rightarrow$ acceder a
        memoria de otros procesos o del kernel.
    \item Deshabilitar las interrupciones $\Rightarrow$ tomar la CPU
        para siempre.
    \item Tocar los registros de los dispositivos $\Rightarrow$
        manipular hardware directamente.
\end{itemize}

Sin un modo privilegiado que el SO use exclusivamente, no hay forma
de \emph{forzar} que un programa pase por las syscalls y los chequeos
correspondientes. Por eso \textbf{toda CPU moderna tiene al menos dos
modos} (user / kernel) y muchas tienen más (x86: 4 rings; ARM: 4
niveles de excepción; RISC-V: M, S, U).

\textbf{Excepciones marginales}: sistemas embebidos con software
totalmente confiable y auditado (un controlador empotrado de un
único programa, p.\ ej.) pueden vivir con un solo modo, pero no son
SOs multi-usuario en sentido estricto.
\end{respuesta}

% =====================================================================
\section{Ejercicio 2 --- ¿Monolítico es más eficiente que microkernel?}
% =====================================================================

\textbf{Enunciado.} Un diseño monolítico ofrece generalmente mayor
eficiencia que un microkernel. ¿Es válida esta afirmación? Justificar.

\subsection*{Análisis}

En un \textbf{kernel monolítico} (Linux clásico, Windows NT,
\texttt{xv6}), todos los subsistemas ---scheduler, gestor de memoria,
FS, drivers, networking--- están compilados en un único binario y
corren en \textbf{modo kernel}, en el mismo espacio de direcciones.
La comunicación entre ellos es por \emph{llamadas a función directa}
y acceso a estructuras globales (tabla de procesos, tabla de archivos
abiertos, etc.).

En un \textbf{microkernel} (Mach, MINIX, L4, QNX), el núcleo se
reduce al mínimo: planificador, IPC, MMU básica. El resto (FS,
drivers, networking, incluso el gestor de procesos) corre como
\textbf{procesos servidores en modo usuario}. Toda operación que
involucre múltiples subsistemas pasa por IPC entre procesos.

\subsection*{Por qué el monolítico es más eficiente}

\begin{itemize}[leftmargin=*]
    \item \textbf{Sin cambios de modo extra}: en monolítico, la
        invocación entre subsistemas del kernel es una llamada a
        función. En microkernel, requiere \textbf{traps adicionales}
        (uno por cada \texttt{send}/\texttt{receive}).
    \item \textbf{Sin cambios de contexto extra}: en monolítico, no
        hay context switch al invocar otro subsistema. En microkernel,
        sí: el control se transfiere al proceso servidor.
    \item \textbf{Sin copia de datos entre espacios}: en monolítico,
        un puntero a un buffer se pasa directamente. En microkernel,
        muchas veces hay que \emph{copiar} los datos al espacio del
        servidor (el espacio de direcciones es distinto).
\end{itemize}

\textbf{Ejemplo numérico}: una syscall \texttt{getpid()} simple
en un monolítico cuesta $\sim 1$ trap (decenas de nanosegundos). En
un microkernel típico: $\sim 2$ traps + $2$ context switches +
copia de mensaje (centenas de nanosegundos, hasta microsegundos en
implementaciones ingenuas).

\subsection*{Lo que el microkernel ofrece a cambio}

A pesar de la menor eficiencia, los microkernels tienen ventajas:

\begin{itemize}[leftmargin=*,nosep]
    \item \textbf{Modularidad}: cada servicio es un proceso aislado.
    \item \textbf{Robustez}: un driver que crashea no tira el sistema
        (el kernel sigue corriendo; se reinicia el servicio).
    \item \textbf{Seguridad}: cada servicio tiene sus propios
        privilegios; un bug en un driver no compromete todo.
    \item \textbf{Portabilidad}: agregar/quitar servicios sin
        recompilar el kernel.
\end{itemize}

\begin{respuesta}
La afirmación es \textbf{válida}: un kernel monolítico es
generalmente más eficiente que un microkernel para la misma carga,
porque evita los cambios de modo, los context switches y la copia
de datos entre espacios de direcciones que requiere el IPC del
microkernel.

A cambio, el microkernel gana en \emph{modularidad, robustez,
seguridad y portabilidad}, lo que justifica su uso en escenarios
donde esos atributos son críticos (sistemas embebidos certificados,
sistemas de tiempo real, sistemas de alta seguridad).
\end{respuesta}

% =====================================================================
\section{Ejercicio 3 --- Instrucciones privilegiadas}
% =====================================================================

\textbf{Enunciado.} Determinar cuáles de las siguientes instrucciones
deberían ser privilegiadas.

\subsection*{Criterio general}

Una instrucción debe ser \textbf{privilegiada} si su ejecución por
parte de un proceso de usuario podría \textbf{comprometer la
seguridad o estabilidad del sistema}. La regla mnemónica: si la
instrucción puede afectar a \emph{otros procesos} o al \emph{kernel},
debe ser privilegiada.

\subsection*{3.a --- Setear el valor del timer}

\begin{respuesta}
\textbf{Privilegiada}. El timer es el corazón del scheduler. Un
proceso que pudiera modificarlo podría:

\begin{itemize}[leftmargin=*,nosep]
    \item Aumentar su quantum efectivo y monopolizar la CPU.
    \item Detener el timer y bloquear el sistema entero.
    \item Falsificar la hora del sistema.
\end{itemize}
\end{respuesta}

\subsection*{3.b --- Leer el valor del timer}

\begin{respuesta}
\textbf{No privilegiada}. Leer la hora actual es una operación
inofensiva. Los procesos normales necesitan saber la hora para
muchas operaciones (timestamps de archivos, mediciones de
rendimiento, etc.). Por eso existen syscalls como
\texttt{gettimeofday}, \texttt{clock\_gettime}, \texttt{time}: son
accesibles desde modo usuario sin riesgo.
\end{respuesta}

\subsection*{3.c --- Ejecutar una instrucción de \texttt{trap}}

\begin{respuesta}
\textbf{No privilegiada}. La instrucción de trap (\texttt{syscall},
\texttt{ecall}, \texttt{svc}) es \emph{exactamente} el mecanismo que
los procesos de usuario usan para invocar al kernel. Si fuera
privilegiada, ningún programa podría hacer syscalls. Es la entrada
canónica a modo kernel, controlada por el hardware: el procesador
salta a una dirección fija definida por el kernel, no a una
arbitraria.
\end{respuesta}

\subsection*{3.d --- Deshabilitar interrupciones}

\begin{respuesta}
\textbf{Privilegiada}. Si un proceso pudiera deshabilitar IRQs:

\begin{itemize}[leftmargin=*,nosep]
    \item El timer no podría interrumpirlo $\Rightarrow$
        monopolizaría la CPU.
    \item Los dispositivos de I/O no podrían notificar al kernel.
    \item No habría preemption $\Rightarrow$ multitarea imposible.
\end{itemize}

Es probablemente la instrucción más peligrosa fuera del control del
kernel.
\end{respuesta}

\subsection*{3.e --- Modificar registros de protección de memoria}

\begin{respuesta}
\textbf{Privilegiada}. Los registros que controlan la MMU (puntero
a la tabla de páginas, base/límite de segmentos) definen qué memoria
ve cada proceso. Modificarlos permitiría:

\begin{itemize}[leftmargin=*,nosep]
    \item Acceder a la memoria de otros procesos.
    \item Acceder al espacio del kernel.
    \item Saltarse cualquier mecanismo de aislamiento.
\end{itemize}

Es la base misma de la separación entre procesos.
\end{respuesta}

\subsection*{3.f --- Acceder a registros de control y estado de dispositivos}

\begin{respuesta}
\textbf{Privilegiada}. Los registros de los controladores de
dispositivo (disco, red, GPU, etc.) viven en regiones de memoria
mapeada por hardware (MMIO) o en el espacio de puertos. El acceso
directo permitiría:

\begin{itemize}[leftmargin=*,nosep]
    \item Leer datos del disco saltándose los permisos del FS.
    \item Inyectar paquetes en la red sin pasar por el firewall del
        SO.
    \item Provocar fallas del hardware.
\end{itemize}

Por eso solo el kernel (y sus drivers) tocan estos registros; los
procesos de usuario los acceden indirectamente vía syscalls
(\texttt{read}, \texttt{write}, \texttt{ioctl}).
\end{respuesta}

\subsection*{Resumen}

\begin{center}
\small
\begin{tabular}{@{}clc@{}}
\toprule
\# & Instrucción                                    & ¿Privilegiada? \\
\midrule
a) & Setear el valor del timer                      & \textbf{Sí}    \\
b) & Leer el valor del timer                        & No             \\
c) & Ejecutar una instrucción de \texttt{trap}      & No             \\
d) & Deshabilitar interrupciones                    & \textbf{Sí}    \\
e) & Modificar registros de protección de memoria  & \textbf{Sí}    \\
f) & Acceder a registros de control de dispositivos & \textbf{Sí}    \\
\bottomrule
\end{tabular}
\end{center}

% =====================================================================
\section{Ejercicio 4 --- Propósito y mecanismo de las syscalls}
% =====================================================================

\textbf{Enunciado.} ¿Cuál es el propósito de una llamada al sistema
y cómo se pueden pasar los parámetros?

\subsection*{Propósito}

Una \textbf{llamada al sistema} (syscall) es la interfaz controlada
por la cual un proceso de usuario solicita un servicio al kernel.

Sirve para que el proceso pueda realizar operaciones que requieren
\textbf{privilegios} (acceso a archivos, red, dispositivos,
manipulación de procesos, etc.) sin tener él mismo esos privilegios:

\begin{itemize}[leftmargin=*,nosep]
    \item El proceso prepara los argumentos en sus registros (o en
        el stack o en memoria).
    \item Ejecuta la instrucción de \emph{trap}.
    \item El hardware cambia a modo supervisor y salta a una
        dirección fija del kernel.
    \item El kernel valida los argumentos, ejecuta la operación en
        nombre del proceso y devuelve el resultado.
    \item El control vuelve al proceso, que ahora está nuevamente
        en modo usuario.
\end{itemize}

La syscall es el \textbf{único punto de entrada controlado} al
kernel. Esto garantiza que (a) los chequeos de permisos siempre se
hagan, (b) los argumentos sean validados, y (c) no se pueda saltar a
código arbitrario del kernel.

\subsection*{Cómo se pasan los parámetros}

Hay tres mecanismos clásicos, todos definidos por la \emph{ABI}
(\emph{Application Binary Interface}) de cada arquitectura:

\subsubsection*{1) En registros (la forma estándar moderna)}

Los argumentos se colocan en registros designados \emph{antes} de la
instrucción de trap, y el número de syscall en otro registro.

\begin{center}
\small
\begin{tabular}{@{}lll@{}}
\toprule
Arquitectura & Registros de args. (en orden) & Registro de syscall \# \\
\midrule
x86-64 (Linux) & \texttt{rdi, rsi, rdx, r10, r8, r9}  & \texttt{rax} \\
ARM (AArch64)  & \texttt{x0--x5}                        & \texttt{x8}  \\
RISC-V         & \texttt{a0--a5}                        & \texttt{a7}  \\
\bottomrule
\end{tabular}
\end{center}

\textbf{Valor de retorno}: usualmente en el mismo registro del
primer argumento (\texttt{rax}, \texttt{x0}, \texttt{a0}). El
proceso lee de ahí el resultado cuando reanuda en modo usuario.

\textbf{Ventajas}: rápido (no toca memoria) y simple para syscalls
con pocos argumentos (hasta 6 normalmente).

\subsubsection*{2) En el stack}

El proceso empuja los argumentos al stack y ejecuta el trap. El
kernel los lee del stack del usuario.

\textbf{Uso histórico}: x86 32-bit (\texttt{int 0x80}). Más lento y
requiere validar punteros al stack del usuario.

\subsubsection*{3) Vía una estructura en memoria}

Para syscalls con muchos argumentos o argumentos grandes (e.g.,
\texttt{ioctl}, \texttt{readv}/\texttt{writev}), se pasa un puntero
a una estructura en el espacio del proceso, y el kernel la lee
después de validar el puntero.

\subsection*{Validación de los argumentos}

\begin{importante}
El kernel \textbf{nunca debe confiar} en los argumentos del proceso.
Antes de usarlos, valida:

\begin{itemize}[leftmargin=*,nosep]
    \item Que los punteros estén en el rango virtual del proceso
        (no apunten al kernel).
    \item Que los rangos de tamaño sean razonables (e.g., un
        \texttt{count} de \texttt{read} no sea negativo o
        excesivo).
    \item Que el proceso tenga permiso para la operación (e.g., el
        FD esté abierto para lectura).
\end{itemize}

La falta de validación es una de las clases más comunes de bugs en
kernel.
\end{importante}

\begin{respuesta}
\textbf{Propósito}: permitir que un proceso de usuario acceda a
servicios privilegiados del kernel (FS, red, dispositivos, gestión
de procesos) sin tener esos privilegios él mismo. La syscall es el
único punto de entrada controlado al kernel; el hardware (vía la
instrucción de trap) garantiza que el control salte a una dirección
fija definida por el SO.

\textbf{Paso de parámetros}: vía registros (lo más común y rápido,
hasta 6 argumentos), vía stack (histórico) o vía estructura en
memoria (para argumentos complejos). El número de syscall se coloca
en un registro convenido por la ABI. El kernel devuelve el
resultado en el mismo registro del primer argumento.
\end{respuesta}

% =====================================================================
\section{Ejercicio 5 --- Shell ejecutando \texttt{cmd1 \& cmd2}}
% =====================================================================

\textbf{Enunciado.} Describir los pasos y llamadas al sistema que
hace un shell en un comando de la forma \texttt{cmd1 \& cmd2}.

\subsection*{Análisis: qué significa \texttt{\&}}

En bash y la mayoría de shells de UNIX, el operador \texttt{\&}
\emph{separa} comandos (similar a \texttt{;}) y marca al comando
de la izquierda para ejecutarse en \textbf{background}:

\begin{itemize}[leftmargin=*,nosep]
    \item \texttt{cmd1 \&}: ejecuta cmd1 en background; el shell no
        espera.
    \item \texttt{cmd1 \& cmd2}: ejecuta cmd1 en background y luego
        cmd2 en foreground (con \texttt{wait}).
\end{itemize}

\begin{importante}
No confundir con \texttt{cmd1 \&\& cmd2} (con dos ampersands), que
es \emph{secuencial-condicional}: ejecuta cmd2 solo si cmd1
terminó con código 0.
\end{importante}

\subsection*{Secuencia de pasos del shell}

\paragraph{1. Parseo de la línea de comando.}
El shell tokeniza la línea y reconoce el operador \texttt{\&} como
separador con semántica de background para cmd1.

\paragraph{2. Ejecución de cmd1 en background.}

\begin{enumerate}[leftmargin=*,nosep,label=2.\arabic*)]
    \item \texttt{fork()}: el shell crea un proceso hijo
        $P_1$ (mismo programa, copia del espacio).
    \item En $P_1$ (hijo): \texttt{exec("cmd1", args)}: reemplaza
        la imagen del proceso con el binario de cmd1.
    \item En el shell (padre): \textbf{NO} hace \texttt{wait()}; el
        shell guarda el PID en su tabla de \emph{jobs}.
    \item El shell imprime algo como \texttt{[1] 12345} (job number
        y PID).
    \item El shell continúa procesando la siguiente parte de la
        línea (cmd2).
\end{enumerate}

\paragraph{3. Ejecución de cmd2 en foreground.}

\begin{enumerate}[leftmargin=*,nosep,label=3.\arabic*)]
    \item \texttt{fork()}: el shell crea otro hijo $P_2$.
    \item En $P_2$: \texttt{exec("cmd2", args)}.
    \item En el shell: \texttt{waitpid($P_2$\_pid, \&status, 0)}:
        bloquea hasta que cmd2 termine.
    \item Cuando $P_2$ termina, el shell vuelve a mostrar el prompt.
\end{enumerate}

\paragraph{4. Cuando cmd1 (background) eventualmente termina:}

\begin{itemize}[leftmargin=*,nosep]
    \item El kernel envía \texttt{SIGCHLD} al shell.
    \item El shell (cuyo handler de \texttt{SIGCHLD} llama a
        \texttt{waitpid} en modo no bloqueante \texttt{WNOHANG})
        recoge el exit status y libera al zombie.
    \item En la próxima impresión del prompt, el shell notifica al
        usuario: \texttt{[1]+ Done   cmd1}.
\end{itemize}

\subsection*{Diagrama del flujo}

\begin{center}
\begin{tikzpicture}[
    every node/.style={font=\ttfamily\small,align=center},
    box/.style={draw,rounded corners,minimum width=2.2cm,minimum height=0.8cm},
    sh/.style={box,fill=c1!25},
    bg/.style={box,fill=c3!25},
    fg/.style={box,fill=c2!25},
    arr/.style={->,thick}
]
\node[sh] (shell)   at (0,4) {Shell};
\node[bg] (p1)      at (-3,2) {$P_1$: cmd1 \\ \scriptsize (background)};
\node[fg] (p2)      at (3,2)  {$P_2$: cmd2 \\ \scriptsize (foreground)};
\node[sh] (shell2)  at (0,0)  {Shell \\ \scriptsize (vuelve prompt)};

\draw[arr] (shell.west) -- node[above,font=\scriptsize]{\texttt{fork}+\texttt{exec}} (p1.north);
\draw[arr] (shell.east) -- node[above,font=\scriptsize]{\texttt{fork}+\texttt{exec}} (p2.north);
\draw[arr,dashed] (p2.south) -- node[right,font=\scriptsize]{\texttt{waitpid}} (shell2.north);
\draw[arr,dashed,gray] (p1.south) to[bend right=30] node[left,font=\scriptsize]{\texttt{SIGCHLD} luego} (shell2.west);
\end{tikzpicture}
\end{center}

\subsection*{Resumen de syscalls usadas}

\begin{center}
\small
\begin{tabular}{@{}llp{7cm}@{}}
\toprule
Syscall    & Quién la invoca   & Función \\
\midrule
\texttt{fork}     & shell                 & crear hijo para cmd1, luego otro para cmd2 \\
\texttt{exec}     & cada hijo             & reemplazar imagen con el binario del comando \\
\texttt{waitpid}  & shell (solo para cmd2) & esperar fin de cmd2 (foreground) \\
\texttt{waitpid} (\texttt{WNOHANG}) & shell en handler de \texttt{SIGCHLD} & recoger exit status de cmd1 cuando termine \\
\bottomrule
\end{tabular}
\end{center}

\begin{respuesta}
El shell ejecuta \texttt{cmd1 \& cmd2} en dos rondas de
\texttt{fork+exec}: la primera, sin \texttt{wait()} (cmd1 en
background, el shell guarda el job y sigue); la segunda, con
\texttt{waitpid} bloqueante (cmd2 en foreground). El zombie de
cmd1 se cosecha más tarde vía el handler de \texttt{SIGCHLD} o en
el próximo prompt.
\end{respuesta}

% =====================================================================
\section{Ejercicio 6 --- ¿\texttt{read} hace busy-wait?}
% =====================================================================

\textbf{Enunciado.} En una llamada al sistema \texttt{read(fd,
buffer, count)} el proceso queda haciendo una espera ocupada hasta
que el kernel finalice la operación. Justifique la respuesta.

\subsection*{Análisis}

\begin{respuesta}
\textbf{Falso}. El proceso \textbf{no} hace busy-wait. Pasa al
estado \texttt{SLEEPING} (bloqueado en una \emph{wait queue}) y
libera la CPU. Mientras espera, otros procesos pueden ejecutar.
\end{respuesta}

\textbf{Por qué no es busy-wait}:

El propósito mismo de un SO multitarea es \emph{no desperdiciar CPU}.
Si \texttt{read} hiciera busy-wait, el proceso quemaría todo su
quantum girando sin hacer trabajo útil, mientras la CPU podría
estar ejecutando algún otro proceso productivo. Sería el peor diseño
posible.

\textbf{Qué hace el kernel realmente}:

\begin{enumerate}[leftmargin=*]
    \item El proceso ejecuta el trap de \texttt{read}.
    \item El kernel verifica si los datos ya están en el
        \emph{buffer cache}:
        \begin{itemize}[nosep]
            \item Si SÍ (cache hit): copia al \texttt{buffer} del
                usuario y retorna inmediatamente.
            \item Si NO (cache miss): sigue.
        \end{itemize}
    \item El kernel envía un comando al controlador del disco para
        leer el bloque correspondiente.
    \item El kernel pone al proceso en una \emph{wait queue}
        asociada a esa operación: \texttt{p->state = SLEEPING},
        encola en \texttt{wq}.
    \item El kernel llama al scheduler ($\to$ \texttt{yield}). Otro
        proceso obtiene la CPU.
    \item Eventualmente, el dispositivo termina la operación y
        dispara una IRQ.
    \item El ISR del driver marca los datos como disponibles y
        llama a \texttt{wakeup(wq)}, que pasa al proceso a
        \texttt{RUNNABLE}.
    \item Cuando el scheduler vuelve a elegirlo, el proceso reanuda
        dentro de \texttt{read}, copia los datos al \texttt{buffer}
        de usuario y retorna.
\end{enumerate}

\begin{idea}
\textbf{Por qué bloquear en vez de gastar CPU}: la I/O de disco es
de órdenes de magnitud más lenta que la CPU. Un acceso a un HDD
toma $\sim 10$\,ms; en ese tiempo la CPU puede ejecutar $\sim 10^7$
instrucciones útiles. Si el proceso girara en busy-wait, todas esas
instrucciones se perderían. El bloqueo es la clave de la
\emph{utilización eficiente} de la CPU.
\end{idea}

% =====================================================================
\section{Ejercicio 7 --- ¿\texttt{wakeup} pasa procesos a RUNNING?}
% =====================================================================

\textbf{Enunciado.} La operación \texttt{wakeup()} pasa procesos
desde el estado \texttt{SLEEPING} a \texttt{RUNNING}. Justificar.

\subsection*{Análisis}

\begin{respuesta}
\textbf{Falso}. \texttt{wakeup()} pasa procesos de \texttt{SLEEPING}
a \texttt{RUNNABLE} (también llamado \texttt{READY}), \textbf{no
directamente a RUNNING}. Para que un proceso pase a \texttt{RUNNING}
hace falta que el \textbf{scheduler} lo seleccione de la cola
\texttt{READY} y le asigne CPU.
\end{respuesta}

\subsection*{Diagrama de estados de un proceso}

\begin{center}
\begin{tikzpicture}[
    every node/.style={font=\ttfamily\small},
    state/.style={draw,rounded corners,minimum width=2.4cm,minimum height=1cm,align=center},
    arr/.style={->,thick}
]
\node[state,fill=c1!25] (run)   at (0,0)   {\textbf{RUNNING}};
\node[state,fill=c2!25] (rdy)   at (-4.5,0) {\textbf{RUNNABLE}};
\node[state,fill=c3!25] (slp)   at (4.5,0)  {\textbf{SLEEPING}};

\draw[arr] (rdy) -- node[above,font=\scriptsize]{scheduler elige} (run);
\draw[arr,bend left] (run.south) to[bend right=25] node[below,font=\scriptsize]{quantum / yield} (rdy.south);

\draw[arr] (run) -- node[above,font=\scriptsize]{\texttt{read}, \texttt{wait}, etc.} (slp);
\draw[arr,bend left,cred] (slp.south) to[bend left=25] node[below,font=\scriptsize]{\textbf{\texttt{wakeup()}}} (rdy.south);
\end{tikzpicture}
\end{center}

Lo crítico: \texttt{wakeup()} NO conecta directamente con
\texttt{RUNNING}. Va a \texttt{RUNNABLE} primero.

\subsection*{Por qué no es directo}

\begin{itemize}[leftmargin=*]
    \item Puede haber muchos procesos \texttt{RUNNABLE}; solo uno
        (o tantos como CPUs haya) puede estar \texttt{RUNNING}
        simultáneamente.
    \item \texttt{wakeup} puede despertar varios procesos a la vez
        (todos los que esperaban en la misma \emph{wait queue}).
        No es posible ponerlos a todos en \texttt{RUNNING} de
        inmediato.
    \item La decisión de quién ejecuta es del scheduler, no del
        que despertó. El scheduler aplica su política (prioridad,
        RR, MLFQ, etc.) para elegir.
\end{itemize}

\subsection*{Excepción: \emph{preemption} al despertar}

Algunos kernels (Linux con \texttt{CONFIG\_PREEMPT}) verifican,
después de \texttt{wakeup}, si el proceso recién despertado tiene
\emph{mayor prioridad} que el actualmente en ejecución. Si es así,
disparan inmediatamente un context switch: el despertado entra a
\texttt{RUNNING} y el actual vuelve a \texttt{RUNNABLE}.

Aún en ese caso, la transición lógica es
\texttt{SLEEPING} $\to$ \texttt{RUNNABLE} $\to$ \texttt{RUNNING} (y
la última depende del scheduler, no del wakeup).

% =====================================================================
\section{Ejercicio 8 --- ¿Qué algoritmos generan starvation?}
% =====================================================================

\textbf{Enunciado.} ¿Cuáles de los siguientes algoritmos de
planificación de uso de CPU pueden generar \emph{starvation}
(\emph{livelocks})?

\begin{enumerate}[leftmargin=*,nosep,label=\alph*)]
    \item FCFS
    \item Shortest Job/Burst First
    \item Round Robin
    \item Prioridades
\end{enumerate}

\subsection*{Análisis: qué es starvation}

\textbf{Starvation}: un proceso espera indefinidamente recibir CPU
mientras otros sí la reciben.

Para que haya starvation hace falta que la política de selección
\emph{discrimine sistemáticamente} contra algún proceso, sin
mecanismo de \emph{aging} que mitigue.

\subsection*{8.a --- FCFS}

\begin{respuesta}
\textbf{No genera starvation}. FCFS es estrictamente por orden de
llegada. Cada proceso espera, como mucho, los $n - 1$ procesos que
llegaron antes (suma de sus ráfagas). En tiempo finito todos ejecutan.

\textbf{Aclaración}: FCFS sufre el \emph{convoy effect} (procesos
cortos esperan detrás de uno largo y aumentan su turnaround
promedio), pero eso es \emph{ineficiencia}, no \emph{starvation}:
nadie espera infinitamente.
\end{respuesta}

\subsection*{8.b --- Shortest Job/Burst First (SJF / SRTF)}

\begin{respuesta}
\textbf{SÍ puede generar starvation}. SJF siempre elige el proceso
con la ráfaga más corta. Si llegan continuamente procesos cortos,
los procesos largos \emph{nunca son seleccionados}: cada vez que
llegan al frente, hay alguien más corto adelante.

Es el caso paradigmático de starvation en libros de SO. La solución
estándar es \emph{aging}: aumentar la prioridad efectiva del proceso
en función del tiempo que lleva esperando.
\end{respuesta}

\subsection*{8.c --- Round Robin}

\begin{respuesta}
\textbf{No genera starvation}. RR es estrictamente justo: cada
proceso obtiene CPU exactamente una vez cada $\leq N \cdot q$
unidades de tiempo (con $N$ procesos en la cola y quantum $q$).
Ningún proceso queda postergado: el peor caso de espera está
acotado.
\end{respuesta}

\subsection*{8.d --- Prioridades}

\begin{respuesta}
\textbf{SÍ puede generar starvation}, si no hay aging. Procesos
de baja prioridad pueden quedar indefinidamente postergados si
constantemente llegan procesos de mayor prioridad.

La solución estándar (también clásica): \emph{aging}, incrementar la
prioridad del proceso a medida que envejece en la cola. Garantiza
que tarde o temprano alcanzará la prioridad necesaria para ejecutar.
\end{respuesta}

\subsection*{Resumen}

\begin{center}
\small
\begin{tabular}{@{}lcl@{}}
\toprule
Algoritmo                 & ¿Starvation? & Razón / Mitigación \\
\midrule
FCFS                      & No           & Estrictamente por orden de llegada \\
SJF / SRTF                & \textbf{Sí}  & Procesos largos postergados; \emph{aging} \\
Round Robin               & No           & Justo por construcción; espera acotada \\
Prioridades (sin aging)   & \textbf{Sí}  & Baja prioridad postergada; \emph{aging} \\
\bottomrule
\end{tabular}
\end{center}

\begin{importante}
\textbf{Starvation vs.\ livelock}. La consigna los menciona juntos
pero son distintos:

\begin{itemize}[leftmargin=*,nosep]
    \item \textbf{Starvation}: el proceso espera CPU/recurso pero
        nunca lo obtiene.
    \item \textbf{Livelock}: los procesos sí ejecutan, pero
        constantemente se ceden el paso o repiten estados sin
        progresar.
\end{itemize}

En algoritmos de planificación, el problema relevante es starvation,
no livelock (los livelocks típicamente aparecen en protocolos de
sincronización mal diseñados).
\end{importante}

% =====================================================================
\section{Ejercicio 9 --- Race condition en sistema bancario}
% =====================================================================

\textbf{Enunciado.} Un sistema bancario usa archivos para representar
el saldo de una cuenta bancaria.

\begin{lstlisting}[style=cstyle]
deposit(account, amount):
    m = read(account)
    write(account, m + amount)

withdraw(account, amount):
    m = read(account)
    if m > amount:
        write(account, m - amount)
\end{lstlisting}

\textbf{(a)} Determinar si existe race condition en un sistema
multitarea preemptivo.

\textbf{(b)} Ofrecer una solución usando algún mecanismo de
sincronización.

\subsection*{9.a --- ¿Hay race condition?}

\begin{respuesta}
\textbf{Sí, hay race condition} de tipo \emph{lost update}. Ambas
operaciones son secuencias \textbf{read-modify-write} no atómicas
sobre el mismo archivo. En un sistema preemptivo, el scheduler
puede interrumpir entre el \texttt{read} y el \texttt{write} de
una tarea para correr la otra, produciendo intercalado nocivo.
\end{respuesta}

\subsection*{Traza concreta del bug}

Saldo inicial \texttt{account} $= 100$.

Sean dos operaciones concurrentes:
\begin{itemize}[leftmargin=*,nosep]
    \item $T_1$: \texttt{deposit(account, 30)} (deposita $30$).
    \item $T_2$: \texttt{withdraw(account, 50)} (retira $50$).
\end{itemize}

Resultado esperado: $100 + 30 - 50 = 80$.

\begin{center}
\small
\begin{tabular}{@{}clll@{}}
\toprule
$t$ & $T_1$ (\texttt{deposit 30})                & $T_2$ (\texttt{withdraw 50}) & saldo \\
\midrule
$t_1$ & \texttt{m = read(account)} $\to$ 100      &                                   & 100 \\
$t_2$ & \multicolumn{2}{l}{\textbf{preempt: scheduler cambia a $T_2$}}                & 100 \\
$t_3$ &                                            & \texttt{m = read(account)} $\to$ 100 & 100 \\
$t_4$ &                                            & \texttt{if m > 50} (verdadero)        & 100 \\
$t_5$ &                                            & \texttt{write(account, 50)}            & \textbf{50} \\
$t_6$ & \multicolumn{2}{l}{\textbf{preempt: scheduler cambia a $T_1$}}                & 50 \\
$t_7$ & \texttt{write(account, 100 + 30 = 130)}    &                                   & \textbf{130} \\
\bottomrule
\end{tabular}
\end{center}

\textbf{Resultado real}: saldo $= 130$. Diferencia con el esperado
($80$): $\$50$ que aparecieron \textbf{de la nada}. Otra traza
posible (a la inversa) llevaría a saldo $= 50$, es decir, $\$30$
\textbf{desaparecen}.

Ambas trazas son posibles y dependen del scheduling. Esto es
exactamente el patrón clásico de \emph{lost update}: las dos
operaciones leen el mismo valor inicial y se pisan al escribir.

\subsection*{9.b --- Solución con un mutex (lock por cuenta)}

\begin{lstlisting}[style=cstyle]
/* Cada cuenta tiene su propio mutex. */
mutex_t account_lock[N_ACCOUNTS];

void deposit(int account, double amount) {
    acquire(&account_lock[account]);
    /* --- region critica --- */
    double m = read(account);
    write(account, m + amount);
    /* --- fin region critica --- */
    release(&account_lock[account]);
}

void withdraw(int account, double amount) {
    acquire(&account_lock[account]);
    /* --- region critica --- */
    double m = read(account);
    if (m > amount) {
        write(account, m - amount);
    }
    /* --- fin region critica --- */
    release(&account_lock[account]);
}
\end{lstlisting}

\textbf{Por qué funciona}: el mutex hace que el bloque
\texttt{read} + (cómputo) + \texttt{write} sea atómico respecto de
cualquier otra invocación sobre \emph{la misma cuenta}. Si dos
threads invocan \texttt{deposit} y \texttt{withdraw} sobre
\texttt{account}, el segundo en llegar espera hasta que el primero
libere el lock; entonces lee el saldo \emph{ya actualizado} y opera
correctamente.

\subsection*{Alternativas más robustas}

\begin{enumerate}[leftmargin=*,nosep]
    \item \textbf{\texttt{flock()} del archivo}: si el saldo está
        \emph{en un archivo}, \texttt{flock(fd, LOCK\_EX)} bloquea
        a otros procesos (no solo threads del mismo proceso) que
        intenten tomarlo. Esencial cuando el saldo es accedido por
        \emph{procesos distintos}.

    \item \textbf{Transferencia entre cuentas}: si hubiera
        \texttt{transfer(from, to)}, hay que tomar \emph{dos}
        locks. Para evitar deadlocks, imponer \textbf{orden global}
        (tomar primero el lock de la cuenta de menor ID).

    \item \textbf{Transacciones (estilo BD)}: para bancos reales,
        las operaciones se encapsulan en transacciones
        \emph{ACID}, gestionadas por un motor de DB. El motor
        garantiza atomicidad, consistencia, aislamiento y
        durabilidad.
\end{enumerate}

\begin{idea}
\textbf{¿Por qué un lock por cuenta y no un lock global?} Si todos
los accesos al banco compitieran por el mismo mutex, dos retiros de
cuentas distintas se serializarían innecesariamente. Con un lock
por cuenta, dos operaciones sobre cuentas distintas \emph{pueden
ejecutarse en paralelo}, mejorando el throughput.
\end{idea}

% =====================================================================
\section{Ejercicio 10 --- Spinlocks en uniproc y deshabilitar IRQs en multiproc}
% =====================================================================

\textbf{Enunciado.} Explicar por qué los spinlocks no son apropiados
en sistemas uni-procesadores y por qué deshabilitar interrupciones
no alcanza en multi-procesadores.

\subsection*{Parte 1: spinlocks en uniprocesador}

\textbf{Recordatorio.} Un spinlock implementa exclusión mutua con
busy-wait:

\begin{lstlisting}[style=cstyle]
void acquire(lock_t *l) {
    while (atomic_swap(l, 1) == 1)
        ; /* spin */
}
void release(lock_t *l) {
    *l = 0;
}
\end{lstlisting}

\textbf{Problema en uniprocesador}: hay \emph{un solo flujo} de
ejecución a la vez. Considerá el siguiente escenario:

\begin{enumerate}[leftmargin=*,nosep]
    \item Thread $A$ toma el lock $l$ y entra a la sección crítica.
    \item El timer interrumpe; el scheduler conmuta a thread $B$.
    \item $B$ ejecuta \texttt{acquire(l)}. Como $l = 1$, entra al
        \texttt{while} y \emph{gira}.
    \item $B$ \textbf{quema todo su quantum} girando sin hacer nada
        útil. La CPU está ocupada $100\%$ pero el progreso es nulo.
\end{enumerate}

\textbf{El lock no se libera hasta que $A$ vuelva a ejecutar}, lo
cual no sucederá mientras $B$ esté en la CPU. Recién al expirar el
quantum de $B$, el scheduler puede planificar a $A$, que termina
la SC y libera el lock.

\textbf{Caso peor}: si $B$ tiene MAYOR prioridad que $A$ y el
scheduler nunca lo desaloja, \emph{$A$ no vuelve a ejecutar nunca}.
$B$ gira para siempre. Es una forma de \textbf{inversión de
prioridad} que produce livelock efectivo.

\textbf{Conclusión}: en uniprocesador conviene usar \emph{sleep
locks} (\texttt{pthread\_mutex\_t}, semáforos): cuando un thread no
puede tomar el lock, se \textbf{bloquea} (estado \texttt{SLEEPING}),
liberando la CPU para que otros (incluido el dueño del lock) puedan
ejecutar.

\subsection*{Parte 2: deshabilitar IRQs en multiprocesador}

\textbf{Razonamiento clásico} en uniprocesador: si deshabilito las
IRQs, ningún timer ni IRQ de dispositivo puede interrumpirme; nadie
me toma la CPU; por lo tanto puedo ejecutar la SC sin que nadie más
acceda al recurso compartido. \textbf{Esto funciona en uniproc.}

\textbf{Problema en multiprocesador}: la instrucción \texttt{cli} (o
equivalente) deshabilita interrupciones \textbf{sólo en la CPU
local}.

\begin{center}
\begin{tikzpicture}[
    every node/.style={font=\ttfamily\small,align=center},
    cpu/.style={draw,rounded corners,minimum width=3.6cm,minimum height=1.4cm}
]
\node[cpu,fill=c1!25] (A) at (0,0) {CPU $A$ \\ \scriptsize IRQs deshabilitadas \\ \scriptsize en SC};
\node[cpu,fill=c2!25] (B) at (5,0) {CPU $B$ \\ \scriptsize IRQs habilitadas \\ \scriptsize ejecutando en paralelo};
\node[draw,minimum width=8cm,minimum height=0.8cm,fill=c3!20] (mem) at (2.5,-2) {Memoria compartida (resource $r$)};

\draw[->,thick] (A.south) -- (mem);
\draw[->,thick,cred] (B.south) -- (mem);
\node[font=\scriptsize,cred] at (3.5,-1.2) {¡acceso simultáneo!};
\end{tikzpicture}
\end{center}

CPU $A$ con sus IRQs deshabilitadas entra a la SC tranquila. Pero
CPU $B$ \emph{sigue ejecutando} en paralelo, con sus IRQs
habilitadas, y puede tocar la misma estructura sin que nada se lo
impida. \textbf{La exclusión mutua se rompe.}

\textbf{Conclusión}: deshabilitar IRQs es necesario solo para
prevenir la \emph{preemption local}, pero no provee \emph{exclusión
entre CPUs}. En multiprocesador hace falta:

\begin{enumerate}[leftmargin=*]
    \item Una instrucción \textbf{atómica} (\texttt{xchg},
        \texttt{cmpxchg}, LL/SC) para coordinar entre las CPUs.
    \item Un \textbf{spinlock} construido sobre esa instrucción
        atómica.
    \item Opcionalmente, deshabilitar IRQs locales \emph{además}
        (por ejemplo, si el lock puede ser tomado desde un handler
        de interrupción de la misma CPU; combinación
        \texttt{spin\_lock\_irqsave} de Linux).
\end{enumerate}

\subsection*{Síntesis}

\begin{center}
\small
\begin{tabular}{@{}p{4.5cm}p{4.2cm}p{4.2cm}@{}}
\toprule
& \textbf{Uniprocesador} & \textbf{Multiprocesador} \\
\midrule
Deshabilitar IRQs & Suficiente para EM (pero solo para regiones cortas) & \textbf{Insuficiente}: otras CPUs siguen activas \\
Spinlock & \textbf{Mal uso}: gasta CPU sin progreso & Apropiado para SC cortas \\
Sleep lock (mutex) & Apropiado: bloquea sin gastar CPU & Apropiado para SC largas \\
\bottomrule
\end{tabular}
\end{center}

\begin{respuesta}
\textbf{Spinlocks en uniproc:} mal idea. Si un thread tiene el lock
y otro lo intenta tomar, el segundo \emph{gira} consumiendo todo
su quantum sin progresar; mientras gira, el dueño del lock no puede
ejecutarse para liberar. Es desperdicio puro. En uniproc se usan
\emph{sleep locks}.

\textbf{Deshabilitar IRQs en multiproc:} insuficiente. Solo afecta a
la CPU local; las demás CPUs siguen ejecutando en paralelo y pueden
acceder a la misma memoria. Se requieren instrucciones \emph{atómicas}
(test-and-set, CAS, LL/SC) y spinlocks construidos sobre ellas.
\end{respuesta}

% =====================================================================
\section*{Resumen del parcial}
\addcontentsline{toc}{section}{Resumen del parcial}
% =====================================================================

\begin{center}
\footnotesize
\begin{tabular}{@{}clp{9cm}@{}}
\toprule
\# & Tema & Respuesta sintética \\
\midrule
1  & Seguridad y privilegios     & (a) problemas: confidencialidad, integridad, DoS, acceso a dispositivos, escalada. (b) No: se necesitan al menos 2 modos \\
2  & Monolítico vs microkernel   & V: monolítico más rápido por evitar IPC; microkernel gana en modularidad, robustez y seguridad \\
3  & Instrucciones privilegiadas & a:Sí, b:No, c:No, d:Sí, e:Sí, f:Sí \\
4  & Propósito de syscalls       & Acceder a servicios privilegiados controlados; params en registros (ABI), stack o estructura \\
5  & Shell \texttt{cmd1 \& cmd2} & \texttt{fork}+\texttt{exec} para cmd1 sin \texttt{wait}; \texttt{fork}+\texttt{exec}+\texttt{waitpid} para cmd2; \texttt{SIGCHLD} para cmd1 al terminar \\
6  & \texttt{read} y busy-wait   & F: el proceso pasa a SLEEPING; el SO planifica otros mientras espera la I/O \\
7  & \texttt{wakeup} y RUNNING   & F: \texttt{wakeup} pasa SLEEPING $\to$ RUNNABLE; el scheduler lo lleva a RUNNING \\
8  & Algoritmos y starvation     & FCFS y RR no; SJF y prioridades sí (mitigable con aging) \\
9  & Race condition bancaria     & (a) Sí, \emph{lost update} clásico; (b) mutex por cuenta sobre el bloque read-modify-write \\
10 & Spinlocks y deshabilitar IRQs & Spinlock en uniproc desperdicia CPU; deshabilitar IRQs en multiproc solo afecta la CPU local. Se necesitan instrucciones atómicas \\
\bottomrule
\end{tabular}
\end{center}

\end{document}
